% ---------------------------------------------------------------
\section{Guarantees and Correctness Properties}
\label{sec:guarantees}
% ---------------------------------------------------------------

The Bailout Protocol's correctness properties are expressed as three
fundamental guarantees---SAFETY, LIVENESS, and ACCOUNTABILITY---each derived
directly from the state machine's structural invariants and the \fact{}
layer's pre-commit enforcement semantics.  These guarantees are not
operational aspirations; they are architectural properties that hold
provided the \fact{} layer operates correctly, the parent-pointer chain
$\alpha(\cdot)$ remains intact, and the human anchor $h(n)$ is reachable.
We prove each guarantee through a small-step argument that traces the
state machine's transition structure, the immutability of the parent pointer,
and the credential binding on resolution events.

% ---------------------------------------------------------------
\subsection{Guarantee 1: SAFETY --- No Autonomous Collapse}
\label{sec:safety}

\begin{definition}[Safety Property (SAFETY)]
\label{def:safety}
The \bxthree{} Framework guarantees that no node $n$ can remain in a state
permitting consequential action without a human principal in the causal
authorization chain.  Equivalently: the system cannot undergo autonomous
collapse---a condition in which an agent continues consequential action
indefinitely without a human principal's knowledge or approval.
\end{definition}

\begin{maintheorem}[SAFETY]
\label{th:safety}
For all nodes $n$ in the spawn tree, for all reachable states $q$ of the
Bailout Protocol state machine $\mathcal{B}$, if $q \in Q_{\text{consequential}}$
then a human principal $h \in H$ lies on the causal authorization chain for
$n$.

\medskip
\noindent Here $Q_{\text{consequential}} = \{\,\texttt{IDLE}\,\}$ is the set of
states in which the \fact{} layer permits consequential action at $n$, and
$H$ is the set of human principals provisioned as anchors in the framework.
\end{maintheorem}

\begin{Proof}[Proof Sketch]
The proof proceeds by structural induction over the state machine's
transition graph.

\textbf{Base.}  The initial state $q_0 = \texttt{IDLE}$ is reachable only
from \texttt{HUMAN\_ACKNOWLEDGED} via the resolution event
$\sigma_{\text{resolve}}(\texttt{ACK})$ issued by the human anchor $h(n)$
(Transition T9).  This requires authentication against the exclusive human
credential $C_{h(n)}$, which no machine actor possesses.  Therefore,
$q_0$ implies human presence on the authorization chain by construction.

\textbf{Inductive step.}  We examine each non-terminal state $q$ and show
that every path leading to a state in $Q_{\text{consequential}}$ must pass
through \texttt{HUMAN\_ACKNOWLEDGED} with a human-issued directive.

\begin{itemize}
\item \texttt{BOUNDED}: Consequential action is blocked by the pre-commit
gate.  The only outgoing transitions are T2 (return to \texttt{IDLE} via
human \texttt{ACK}) and T3 (forced timeout, proceeds to
\texttt{BAIL\_PENDING}).  No transition from \texttt{BOUNDED} directly permits
consequential action without human resolution.

\item \texttt{BAIL\_PENDING}: Consequential action is blocked.  The sole
transition T4 is timeout-forced to \texttt{ESCALATING}.  No machine actor can
suppress or redirect this transition (SP-2, Section~\ref{sec:bp-state-machine}).

\item \texttt{ESCALATING}: Consequential action is blocked.  The machine
exits this state only via T6 (propagation reaches anchor node, transitions
to \texttt{HUMAN\_ACKNOWLEDGED}) or T7 (pathway exhausted, terminal failure).
Neither exit path permits consequential action without human acknowledgment.

\item \texttt{HUMAN\_ACKNOWLEDGED}: Consequential action is blocked.  The
machine exits via T8 (human issues binding directive) or T10 (human timeout
expires, terminal failure).  Only T8 produces a non-terminal accepting state.
\end{itemize}

\textbf{Consequences.}  The only non-terminal accepting state is
\texttt{RESOLVED}, reachable exclusively via T8 from
\texttt{HUMAN\_ACKNOWLEDGED} with a $\sigma_{\text{resolve}}$ event
authenticated against $C_{h(n)}$.  \texttt{IDLE} (the state permitting
consequential action) is reachable only by returning through this same human
credentialed path.  No cycle exists in the transition graph that permits a
machine actor to re-authorize consequential action after the state machine
has exited \texttt{IDLE} without human intervention.  $\square$
\end{Proof}

The SAFETY guarantee is structurally enforced by the non-regression property
(SP-1) and the human-only resolution property (SP-3) established in
Section~\ref{sec:bp-state-machine}.  It does not depend on any machine actor
choosing to comply with the protocol; it follows from the deterministic
transition relation alone.

% ---------------------------------------------------------------
\subsection{Guarantee 2: LIVENESS --- Eventual Human Contact}
\label{sec:liveness}

\begin{definition}[Liveness Property (LIVENESS)]
\label{def:liveness}
The \bxthree{} Framework guarantees that every exception state $e$ that
triggers the Bailout Protocol at a node $n$ reaches the designated human
anchor $h(n)$ in finite time.  Formally: for all nodes $n$ and all exception
states $e$ satisfying the trigger condition $\text{TC}(n, x)$, the escalation
path $C(n)$ generated by Algorithm~\ref{alg:escalation} terminates at $h(n)$
within a bounded number of state machine transitions.
\end{definition}

\begin{maintheorem}[LIVENESS]
\label{th:liveness}
For all nodes $n$ and all exception states $e$ at $n$, the Bailout Protocol
state machine reaches the state \texttt{HUMAN\_ACKNOWLEDGED} within at most
$T_{\max}$ transitions, where $T_{\max} = 4 + (d(n) - d(h(n)))$ and
$d(\cdot)$ denotes depth in the spawn tree.

\medskip
\noindent Equivalently: the human anchor $h(n)$ is contacted in finite time for
every genuine exception, regardless of the state or reasoning strategy of any
machine actor on the escalation path.
\end{maintheorem}

\begin{Proof}[Proof Sketch]
The proof proceeds by analyzing the state machine's transition sequence from
the trigger condition firing to human acknowledgment.

\begin{itemize}
\item \textbf{Trigger to escalation (bounds timeout).}  When TC fires at $n$,
the state machine transitions IDLE $\rightarrow$ BOUNDED (T1) and $\tau_{\text{bounds}}$
is armed.  By SP-2 (timeout-forced progression), the machine transitions to
\texttt{BAIL\_PENDING} (T3) when $\tau_{\text{bounds}}$ expires, regardless of
any machine actor's assessment.  This is the first forced transition; it occurs
within at most $T_{\text{bounds}}$ time units.

\item \textbf{Escalation initiation (propagation timeout).}  From
\texttt{BAIL\_PENDING}, the sole transition is to \texttt{ESCALATING} (T4),
forced by $\tau_{\text{prop}}$ timeout expiration.  No machine actor can suppress
this transition.  The state enters \texttt{ESCALATING} within at most
$T_{\text{prop}}$ additional time units.

\item \textbf{Propagation along the parent chain.}  The escalation loop
(Algorithm~\ref{alg:escalation}, T5 self-loop) advances the propagation path
one hop per successful \texttt{SendToParent} invocation.  Each hop moves from a
node $v$ to its parent $\alpha(v)$, strictly decreasing depth:
$d(\alpha(v)) = d(v) - 1$ for all $v$ with $v \neq h(n)$.  The chain terminates
at $h(n)$ where $d(h(n))$ is minimal.  The maximum number of hops is
$d(n) - d(h(n)) \leq D_{\max}$, a finite structural bound established at root
initialization.  Each hop is bounded by $\tau_{\text{prop}}$ timeout; retries
do not extend the total hop count beyond $D_{\max}$.

\item \textbf{Human acknowledgment.}  When propagation reaches the anchor node,
the state machine transitions to \texttt{HUMAN\_ACKNOWLEDGED} (T6).  This
completes the escalation path.
\end{itemize}

\textbf{Combined bound.}  The total number of transitions from the first
trigger to \texttt{HUMAN\_ACKNOWLEDGED} is at most:
\begin{equation}
T_{\max} \leq 1 \;(\text{T1}) \;+\; 1 \;(\text{T3}) \;+\; 1 \;(\text{T4}) \;+\;
\bigl(d(n) - d(h(n))\bigr) \;(\text{T5 loops}) \;+\; 1 \;(\text{T6})
= 4 + \bigl(d(n) - d(h(n))\bigr).
\end{equation}
Since $d(n) - d(h(n))$ is bounded by $D_{\max}$, a finite constant,
$T_{\max}$ is finite for all nodes.  $\square$
\end{Proof}

The LIVENESS guarantee depends on two structural preconditions: the parent
pointer chain $\alpha(\cdot)$ must be computationally intact (each hop resolves
to a valid parent node), and the spawn tree depth must be bounded by $D_{\max}$.
Both are established at node initialization and enforced by the \fact{} layer's
pre-commit hook.  The guarantee does not require that machine actors cooperate
voluntarily; propagation is structurally compelled at each hop by the
pre-commit enforcement at the \fact{} layer.

% ---------------------------------------------------------------
\subsection{Guarantee 3: ACCOUNTABILITY --- Human Always Identifiable}
\label{sec:accountability}

\begin{definition}[Accountability Property (ACCOUNTABILITY)]
\label{def:accountability}
The \bxthree{} Framework guarantees that every exception state that reaches
the Bailout Protocol produces a Forensic Ledger entry that (i) uniquely
identifies the human principal $h(n)$ responsible for the exception's resolution,
(ii) records the complete causal chain from the triggering event to the human
directive, and (iii) is cryptographically sealed against post-hoc modification.
The human principal who authorized the resolution is attributable in the
framework's immutable record for all time.
\end{definition}

\begin{maintheorem}[ACCOUNTABILITY]
\label{th:accountability}
For every exception state $e$ that triggers the Bailout Protocol at a node
$n$, the Forensic Ledger entry chain generated by the protocol satisfies:
\begin{enumerate}
\item[(i)] the human anchor $h(n)$ is uniquely determined by the spawn tree
structure and is recorded in the \texttt{bailout-trigger} entry at the
originating node;
\item[(ii)] every state transition in the protocol's execution is recorded
in the Forensic Ledger before the transition commits (Invariant I4, Section~
\ref{sec:bailout-fact-layer}), forming a complete causal chain; and
\item[(iii)] the entry chain is SHA-256 hash-chained, making any post-hoc
modification of any entry detectable by hash verification.
\end{enumerate}
Consequently, for every exception that entered the Bailout Protocol, the human
principal responsible is uniquely identifiable from the immutable record,
and the complete propagation history is attributable.
\end{maintheorem}

\begin{Proof}[Proof Sketch]
We prove each clause separately.

\medskip
\noindent\textbf{Proof of (i): Unique human anchor identification.}
The human anchor $h(n)$ is provisioned at node $n$'s initialization and stored
in \fact{} layer-managed memory.  It is immutable: the \fact{} layer's
pre-commit hook rejects any runtime attempt to modify $h(n)$ (Invariant I3).
The \texttt{bailout-trigger} Forensic Ledger entry records $h(n)$ explicitly
as part of the entry's structured payload, along with the originating node
$n$ and the trigger condition sub-components.  Because $h(n)$ is established
once and never modified, it is unique and permanent for the node's lifetime.
No machine actor can substitute a different human principal, because no
machine actor at any plane possesses the credential provisioning authority
required to designate or change $h(n)$.  $\square$

\medskip
\noindent\textbf{Proof of (ii): Complete causal chain.}
Invariant I4 (Section~\ref{sec:bailout-fact-layer}) enforces that every state
transition in $\mathcal{B}$ produces at least one corresponding Forensic
Ledger entry \emph{before} the transition commits.  This is enforced by the
\fact{} layer's pre-commit hook: any state transition attempted without a
preceding ledger entry is rejected.  The entry types generated by the protocol
(\texttt{bailout-trigger}, \texttt{timeout-bounds-expired},
\texttt{escalation-sent}, \texttt{human-ack}, \texttt{human-directive})
together form a causally ordered chain: each entry references the node and
state at the time of writing, and propagation entries reference the previous
hop's pathway identifier.  The causal ordering is guaranteed by the
pre-commit enforcement timing: the entry is written and hashed before the
state transition commits, so the record of state $q$ is causally prior to
state $q'$.  $\square$

\medskip
\noindent\textbf{Proof of (iii): Cryptographic sealing.}
Each Forensic Ledger entry $E_i$ contains $\hash{E_{i-1}}$, the SHA-256 hash
of the preceding entry $E_{i-1}$, forming the chain
$E_0, E_1, E_2, \ldots, E_k$.  The genesis entry $E_0$ contains its own hash
as a self-sealing bootstrap.  A post-hoc modification of any entry $E_j$
changes $\hash{E_j}$, which breaks the hash reference in $E_{j+1}$ and all
subsequent entries.  Verification is performed by recomputing the chain hash
forward from $E_0$ and comparing against the stored hash in $E_k$; any
discrepancy indicates tampering.  The hash function is pre-image resistant
and collision resistant under standard cryptographic assumptions, making it
computationally infeasible for a machine actor to modify an entry without
detection.  $\square$
\end{Proof}

The ACCOUNTABILITY guarantee is the conjunction of the two preceding guarantees
plus the \fact{} layer's integrity mechanisms.  SAFETY ensures the human is
present on the authorization chain; LIVENESS ensures the human is contacted;
ACCOUNTABILITY ensures the contact is recorded, attributable, and tamper-evident.
Together, they establish that the \bxthree{} Framework's upstream accountability
guarantee is not merely a design principle but a formally evidenced property of
the runtime execution trace.

% ---------------------------------------------------------------
\subsection{Combined Correctness Statement}
\label{sec:combined-correctness}

The three guarantees are not independent; they form a combined correctness
statement for the Bailout Protocol.

\begin{maintheorem}[Combined Correctness]
\label{th:combined-correctness}
For all nodes $n$ in the spawn tree and all exception states $e$ at $n$:

\begin{enumerate}
\item[SAFETY:] the system cannot execute consequential action without a human
principal on the authorization chain;
\item[LIVENESS:] the exception reaches the designated human anchor $h(n)$ in
finite time; and
\item[ACCOUNTABILITY:] the complete propagation path and the human principal's
resulting directive are recorded in an immutable, cryptographically sealed
Forensic Ledger entry chain.
\end{enumerate}
\end{maintheorem}

The combined statement follows directly from Theorems~\ref{th:safety},
\ref{th:liveness}, and~\ref{th:accountability}.  It establishes that the
upstream accountability guarantee is not an operational convention but an
architecturally enforced property that holds unconditionally within the
framework's stated assumptions.

\begin{conditions-list}[Model Assumptions]
\item The \fact{} layer's pre-commit enforcement operates correctly: the
state machine transitions are evaluated and enforced as specified.
\item The parent-pointer chain $\alpha(\cdot)$ is computationally intact:
each node in the chain resolves to a valid parent node up to the human anchor.
\item The human anchor $h(n)$ is provisioned and reachable at node
initialization.
\item The SHA-256 hash function remains pre-image and collision resistant.
\end{conditions-list}

The assumptions are structural; violations represent failures of the
enforcement substrate rather than failures of the protocol itself.  The
Forensic Ledger records evidence of any assumption violation, enabling
detection and external auditing.  The research agenda motivated by relaxing
these assumptions---particularly the investigation of Byzantine fault
tolerance in the propagation chain and distributed human anchor pools---is
discussed in Section~\ref{sec:bailout-environmental}.